%
\subsection*{Ecuación lineal y respuesta a una entrada}
\label{CH02-B01}
En un intervalo donde pueda escribirse $y'+p(t)y=q(t)$, con $p,q$
continuas, el factor integrante convierte el miembro izquierdo en una
derivada de producto. Con
$P(t)=\int_{t_0}^{t}p(s)\,ds$, el PVI tiene la representación
\[
 y(t)=e^{-P(t)}\left[y_0+\int_{t_0}^{t}e^{P(s)}q(s)\,ds\right].
\]
Así se separan la contribución del dato inicial y la de la entrada, aunque
no haya primitivas elementales. La regla del producto y el teorema
fundamental del cálculo permiten verificarla. Si antes se ha dividido por
el coeficiente de $y'$, sus ceros se estudian en la ecuación original;
el factor integrante no amplía su dominio.
Consultas: soluciones de \exref{M2-E0196} y \exref{M2-E0232}.

%
\subsection*{Exactitud, potencial y factor integrante}
\label{CH02-B02}
Al derivar parcialmente se mantiene fija la otra variable. El diferencial
de $\Phi(x,y)$ es $d\Phi=\Phi_x\,dx+\Phi_y\,dy$. Para
$M\,dx+N\,dy=0$, con $M,N$ de clase $C^1$ en un rectángulo, la igualdad
$M_y=N_x$ permite buscar $\Phi_x=M$, $\Phi_y=N$ y escribir
$\Phi=C$. Al integrar respecto de $x$, la constante es una función de $y$.
Si $\Phi_y\ne0$, el nivel define localmente un gráfico $y(x)$.

Donde los cocientes estén definidos, se pueden ensayar
\[
 \frac{\mu'(x)}{\mu(x)}=\frac{M_y-N_x}{N},\qquad
 \frac{\mu'(y)}{\mu(y)}=\frac{N_x-M_y}{M},
\]
si cada miembro derecho depende sólo de la variable indicada.
Para un factor $\mu(z(x,y))$, la regla de la cadena exige
$\mu'(z)(M z_y-N z_x)=\mu(z)(N_x-M_y)$.
Trabájese donde el factor esté definido y sea no nulo; revísense los
valores excluidos. Un ángulo definido por arctangente requiere una rama
local. La exactitud local no prueba por sí sola exactitud global en un
dominio con huecos. Consulta: solución de \exref{M2-E0021}.

%
\subsection*{Sustituciones y soluciones recuperadas}
\label{CH02-B03}
Derívese primero la variable nueva. Si $z=ax+by+c$, entonces
$z'=a+by'$; si $y=xv$, entonces $y'=v+xv'$ y se trabaja con $x\ne0$.
La ecuación de Bernoulli
\[
 y'+py=qy^n,\qquad z=y^{1-n},\qquad
 z'+(1-n)pz=(1-n)q
\]
se reduce así cuando $n\ne0,1$ y las potencias y divisiones son válidas.
Los casos $n=0,1$ son lineales. Al invertir el cambio, deben recuperarse
signos, ramas y posibles soluciones nulas de la ecuación original.
Una raíz continua no garantiza una solución derivable.

En trayectorias ortogonales, dos pendientes finitas satisfacen
$m_1m_2=-1$; una tangente horizontal corresponde a una vertical, que no se
describe mediante una pendiente finita. Recuérdense la regla de la cadena
y el dominio de las potencias reales. Consultas: soluciones de
\exref{M2-E0019} y \exref{M2-E0207}.

%
\subsection*{Balances por fases y parámetros}
\label{CH02-B04}
En una mezcla homogénea, cantidad de soluto y concentración cumplen
$Q=Vc$. Por tanto, $Q'=Vc'+V'c$: sólo con volumen constante puede
sustituirse $Q'$ por $Vc'$. Un caudal volumétrico multiplicado por una
concentración es una tasa de transporte de soluto. En balances eléctricos
o energéticos, utilícese la ley suministrada y compruébense las dimensiones
de cada constante.

Si una operación cambia en $t=a$ sin aporte instantáneo, el estado
almacenado es continuo y su valor final inicia la fase siguiente; su
derivada puede saltar. Distinguir las fases evita imponer condiciones
incompatibles. Para identificar parámetros hacen falta datos
independientes suficientes. Tratar un reactivo como constante exige que
su variación sea despreciable en el intervalo estudiado, no sólo que
inicialmente sea abundante. Consulta: solución de \exref{M2-E0030}.
